\section*{Acknowledgements}

The constructive framework is that of Osterwalder, Schrader and Seiler, the
curvature criterion that of Bakry and \'Emery, and the uniqueness criterion
that of Dobrushin. The maximal correlation of Gaussian variables is
Gebelein's, the finite-size criterion follows Knabe and Gosset--Mozgunov, the
two-dimensional analysis rests on the heat kernel of the continuum measure of
Driver, Gross, King, Sengupta and L\'evy, and the twisted boundary conditions
are 't~Hooft's; the transfer matrix of a cylinder is built as in L\"uscher's
construction, and the positivity argument for its leading eigenvector is
Jentzsch's. The Lean~4 formalisation of the strong-coupling analysis
described in the introduction is the work of I.~J.~Sessford, and Lean files
for the lattice part of the project were written by J.~M.~Lockwood.

\section*{Author contributions}

Deep Bhattacharjee, who leads the project, posed the problem in the form
treated here, carried out the research and wrote the paper. The proofs are
his: the lattice gap and its constants, the cube condition and Theorems~B
to~E, the transfer matrix of the cylinder, the reduction to one scale, the
ascent and Theorem~F, and the appendices. He made all the computations: the
Lean~4 project of Table~\ref{tab:lean}, the Julia script of
Table~\ref{tab:julia} with its exact rational evaluations, the
numerical checks of the constants, and the figures; and he brought the work
to its present state. James Montgomery Lockwood wrote Lean developments for
the lattice part and proposed the block-projection metrics examined in
earlier versions. Tristan Romanov contributed to the strong-coupling
analysis, in particular the layer bound and the block responses, ran
numerical experiments on the cube gap of $\SU(2)$, and supplied computing
tools. Ikailo John Sessford formalised the strong-coupling analysis on
periodic lattices in Lean~4, in the separate development described in the
introduction. Ushashi Bhattacharya read the manuscript and verified its
arguments and references. Sanjeevan Singha Roy read the manuscript.

\section*{Open Gluon and Anaxia}

Two tools built by D.~B. stand behind the way this project was organised.
The workspace formerly named Open Gluon is now maintained as Anaxia at
\url{https://github.com/anaxia-research/anaxia}. It keeps pinned
copies of the mathematical libraries, the source record of the project and
the checks in one reproducible place. Anaxia, the engine of the eBLLM
project, is his construction engine: it stores each construction as a route
with its input, its output, the obligations that must be discharged and a
test that makes the route fail when an obligation is skipped, and it
certifies finite statements in exact rational arithmetic. He wrote both and
developed them from the constructions of his earlier papers. Three of their
constructions were used here. The first is exact certification: the Julia script recomputes the
constants of this paper, in exact rational arithmetic wherever the value is
rational. The second is the transfer of a boundary response
between a block law and its restriction to the complement of a rare event,
which was developed inside this project in its earlier versions, where the
cube condition of Section~\ref{sec:cube} originated. The third is tightness with observables that do not
escape to infinite energy, the route whose failure test, that a unit
normalisation does not prevent the mass from escaping, is illustrated by
Theorem~\ref{thm:haarlimit}. The engine proposes routes and tests them;
every statement of this paper is proved in the text.

\section*{Declarations}

\emph{Data and code availability.} The supporting Lean~4, Julia and Python
source, manuscript source and verification records are available at
\url{https://github.com/anaxia-research/4Y-M}.
Anaxia, the construction workspace formerly named Open Gluon, is available at
\url{https://github.com/anaxia-research/anaxia}.

\emph{Funding.} No funding of any kind, internal or external, was received
for this work. The authors met its costs themselves; U.~B. and S.~S.~R. took
part in their own time and with their own means.

\emph{Competing interests.} None of the authors has a competing interest,
and neither do the institutions with which they are or were affiliated,
which had no part in the work.

\emph{Ethics.} The work raises no ethical question; it involves no human
participants, no animals and no personal data.
